\documentclass[12pt,a4paper]{article}

% -------------------------------------------------
% Sprache / Kodierung
% -------------------------------------------------
\usepackage[T1]{fontenc}
\usepackage[utf8]{inputenc}
\usepackage[ngerman]{babel}
\usepackage{lmodern}
\usepackage{microtype}

% -------------------------------------------------
% Mathematik
% -------------------------------------------------
\usepackage{amsmath,amssymb,amsthm,mathtools}

% -------------------------------------------------
% Layout / Grafik / Farben
% -------------------------------------------------
\usepackage[a4paper,margin=2.3cm]{geometry}
\usepackage{booktabs}
\usepackage{xcolor}
\usepackage[hidelinks,unicode]{hyperref}
\pdfstringdefDisableCommands{%
  \def\ü{ue}\def\ö{oe}\def\ä{ae}\def\Ü{Ue}\def\Ö{Oe}\def\Ä{Ae}\def\ss{ss}%
}
\usepackage[most]{tcolorbox}
\usepackage{enumitem}
\usepackage{array}

% -------------------------------------------------
% Farben (konsistent mit Vorgängerdokumenten)
% -------------------------------------------------
\definecolor{lückenrot}{RGB}{180,40,40}
\definecolor{bewiesengrün}{RGB}{30,100,50}
\definecolor{warnorange}{RGB}{200,100,10}
\definecolor{hintergrundblau}{RGB}{235,242,250}
\definecolor{hintergrundgelb}{RGB}{255,252,225}
\definecolor{adischviolett}{RGB}{90,20,140}
\definecolor{fehlerrot}{RGB}{200,0,0}
\definecolor{konjviolett}{RGB}{80,0,130}
\definecolor{e8grau}{RGB}{55,55,90}
\definecolor{kingmanblau}{RGB}{20,60,130}
\definecolor{zirkelrot}{RGB}{160,0,0}

% -------------------------------------------------
% tcolorbox-Stile
% -------------------------------------------------
\tcbset{
  lücke/.style={colback=red!8, colframe=lückenrot, fonttitle=\bfseries,
    title={Offene Lücke}},
  ergebnis/.style={colback=green!7, colframe=bewiesengrün,
    fonttitle=\bfseries, title={Bewiesenes Ergebnis}},
  warnung/.style={colback=orange!10, colframe=warnorange,
    fonttitle=\bfseries, title={Achtung / Heuristik}},
  adisch/.style={colback=violet!6, colframe=adischviolett,
    fonttitle=\bfseries, title={2-adische Perspektive}},
  kernidee/.style={colback=hintergrundblau, colframe=blue!60!black,
    fonttitle=\bfseries, title={Kernidee}},
  konj/.style={colback=violet!5, colframe=konjviolett, fonttitle=\bfseries,
    title={Konjektur (unbewiesen)}},
  lte/.style={colback=yellow!8, colframe=e8grau, fonttitle=\bfseries,
    title={LTE-Analyse}},
  hauptsatz/.style={colback=green!10, colframe=bewiesengrün!80!black,
    fonttitle=\bfseries\large, title={Satz (bewiesen)}},
  negativ/.style={colback=red!5, colframe=lückenrot!80!black,
    fonttitle=\bfseries, title={Negatives Ergebnis (bewiesen)}},
  kingman/.style={colback=blue!5, colframe=kingmanblau,
    fonttitle=\bfseries, title={Kingman-Perspektive}},
  tao/.style={colback=teal!5, colframe=teal!60!black,
    fonttitle=\bfseries, title={Verbindung zu Tao (2019)}},
  rpf/.style={colback=cyan!5, colframe=cyan!50!black,
    fonttitle=\bfseries, title={RPF / Spektraltheorie}},
  zirkel/.style={colback=red!12, colframe=zirkelrot,
    fonttitle=\bfseries\color{white}, coltitle=white,
    title={Zirkularit\"{a}tswarnung --- Schritt C ist schwach}},
  bewmat/.style={colback=gray!5, colframe=gray!60!black,
    fonttitle=\bfseries, title={Bewertungsmatrix}},
}

% -------------------------------------------------
% Theorem-Umgebungen
% -------------------------------------------------
\newtheorem{definition}{Definition}[section]
\newtheorem{proposition}[definition]{Proposition}
\newtheorem{theorem}[definition]{Theorem}
\newtheorem{lemma}[definition]{Lemma}
\newtheorem{korollar}[definition]{Korollar}
\newtheorem{bemerkung}[definition]{Bemerkung}
\newtheorem{hypothese}[definition]{Hypothese}
\newtheorem{satz}[definition]{Satz}

\newenvironment{beweis}{\begin{proof}[Beweis]}{\end{proof}}

% -------------------------------------------------
% Makros
% -------------------------------------------------
\newcommand{\vii}{\nu_2}
\newcommand{\U}{\mathcal{U}}
\newcommand{\N}{\mathbb{N}}
\newcommand{\Z}{\mathbb{Z}}
\newcommand{\R}{\mathbb{R}}
\newcommand{\H}{\mathbb{H}}
\newcommand{\Ztwo}{\mathbb{Z}_2}
\newcommand{\abs}[1]{\lvert#1\rvert}
\newcommand{\atwo}[1]{\lvert#1\rvert_2}
\newcommand{\Epi}{\mathbb{E}_\pi}
\newcommand{\TV}{\mathrm{TV}}
\newcommand{\Lam}{\Lambda}
\newcommand{\col}{\mathrm{Col}}

% -------------------------------------------------
% Dokument
% -------------------------------------------------
\title{%
  Das Uniformitätsproblem bei Collatz\\[0.3em]
  \large RPF-Mischung, mod-12-Primitivität und die Grenzen des
  Quaternionen-Arguments\\[0.5em]
  \normalsize Ergänzung zu \textit{collatz\_kingman.tex} und
  \textit{collatz\_lte\_dichte.tex}%
}
\author{Thomas Hoffbauer}
\date{16.\ Juni 2026}

\begin{document}
\maketitle

\begin{abstract}
Taos Ergebnis (2022) zeigt: fast alle $n \in \N$ haben Collatz-Bahnen,
die beliebig klein werden --- im Sinne der logarithmischen Dichte.
Die verbleibende Lücke zur vollen Vermutung ist das
\emph{Uniformitätsproblem}: Maßargumente liefern nur Aussagen für
``fast alle $n$'', aber nicht für alle $n \in \N$.

Wir untersuchen drei Strategien, diese Lücke zu schließen:
\textbf{(A)} LTE-Worst-Cases als ``ewige Ausnahmen'' ausschließen via
Ergodentheorie;
\textbf{(B)} mod-12-Primitivität und exponentiellen Mischungssatz
des Transfer-Operators;
\textbf{(C)} Quaternionen-Norm-Identität als diophantische Schranke.

\textbf{Ergebnis der Analyse:}
Schritt A ist konzeptionell stark, setzt aber Ergodizität voraus.
Schritt B ist das mathematisch solideste Argument und quantifizierbar.
Schritt C enthält eine versteckte Zirkularität und liefert
keinen echten Widerspruch --- er wird in \S\,4 mit roter Warnung markiert.
\end{abstract}

\tableofcontents

\bigskip
\begin{tcolorbox}[warnung, title={Einordnung dieses Dokuments}]
Dieses Dokument enthält eine Kombination aus \textbf{bewiesenen Sätzen}
der Ergodentheorie (grün), \textbf{Anwendungsanalysen} auf Collatz
(orange, teils offen) und \textbf{einer kritischen Fehleranalyse} (rot).
Die Collatz-Vermutung bleibt unbewiesen.
Insbesondere schließt kein Argument in diesem Dokument die volle Lücke.
\end{tcolorbox}

% ====================================================
\section{Die Ergodizitätslücke: präzise Formulierung}
\label{sec:luecke}
% ====================================================

\subsection{Was Tao beweist --- und was nicht}

\begin{tcolorbox}[tao, title={Taos Hauptergebnis (Forum Math.\ Pi, 2022)}]
Für jede Funktion $f: \N \to \R$ mit $f(n) \to \infty$ gilt:
\[
  \mathrm{Dens}_{\log}\bigl(\{n \in \N : \min_{k\geq 0}\col^k(n) \leq f(n)\}\bigr) = 1,
\]
wobei $\mathrm{Dens}_{\log}$ die logarithmische Dichte bezeichnet.
\end{tcolorbox}

Taos Methode arbeitet mit Fourier-Analyse auf $\Ztwo$ und
exponential-sum-Abschätzungen. Im Kern zeigt sie, dass das
Haar-Maß $\mu_H$ auf $\Ztwo^*$ unter der Collatz-Dynamik
``quasi-ergodisch'' ist: Fast alle Startpunkte im 2-adischen Sinne
sind konvergent.

\subsection{Die Ergodizitätslücke}

\begin{tcolorbox}[lücke, title={Das Maß-$0$-Problem (zentral)}]
Sei $E \subseteq \N$ die (hypothetische) Menge aller $n \in \N$,
deren Collatz-Bahn nicht gegen $1$ konvergiert. Taos Ergebnis zeigt:
\[
  \mu_H(E) = 0
  \quad\text{(E hat 2-adisches Haar-Maß Null).}
\]
Aber: $\N \subset \Ztwo^*$ ist eine \emph{abzählbare} Menge, also
selbst eine $\mu_H$-Nullmenge. Deshalb ist $\mu_H(E) = 0$ kein
Widerspruch zu $E \neq \emptyset$.

\medskip
\textbf{Die Lücke:} Man muss zeigen, dass jedes einzelne $n \in \N$
``typisch'' im ergodischen Sinne ist --- es darf auf keiner
$T$-invarianten Nullmenge liegen. Das ist das
\emph{Uniformitätsproblem}.
\end{tcolorbox}

Formal: Birkhoff und Kingman liefern Konvergenz für
$\mu$-\emph{fast alle} Startpunkte. Die Collatz-Vermutung fordert
Konvergenz für \emph{alle} $n \in \N$. Da $\N$ eine $\mu_H$-Nullmenge
ist, sind diese Aussagen grundsätzlich verschieden.

\subsection{Drei Strategien, die Lücke zu schließen}

\begin{center}
\renewcommand{\arraystretch}{1.4}
\begin{tabular}{>{\raggedright}p{1.8cm}
                >{\raggedright}p{5.5cm}
                >{\centering}p{2.0cm}
                >{\raggedright\arraybackslash}p{3.8cm}}
\toprule
\textbf{Schritt} & \textbf{Idee} & \textbf{Status} & \textbf{Bewertung} \\
\midrule
A & LTE-Worst-Cases können nicht ewig dominieren (Ergodentheorie)
  & \textcolor{warnorange}{Bedingt}
  & Setzt Ergodizität voraus \\[2pt]
B & mod-12-Primitivität erzwingt exponentielle Mischung
  & \textcolor{bewiesengrün}{Partiell}
  & Quantifizierbar; formal offen \\[2pt]
C & Quaternionen-Norm als diophantische Schranke
  & \textcolor{lückenrot}{Schwach}
  & Enthält Zirkularität \\
\bottomrule
\end{tabular}
\end{center}

% ====================================================
\section{Schritt A: Lokale Expansion und ``ewige Worst-Cases''}
\label{sec:schrittA}
% ====================================================

\subsection{Die Idee}

Ein hypothetisches $n \in E$ (Ausnahmeset) müsste eine Collatz-Trajektorie
haben, die sich \emph{systematisch und dauerhaft} in den LTE-Worst-Case-Zuständen
aufhält. Der Transfer-Operator würde dann nie zum RPF-stationären Maß $\pi$
konvergieren. Ist das möglich?

\subsection{Formale Formulierung}

Sei $\mathcal{L}: \R^S \to \R^S$ der Transfer-Operator der mod-12-Markov-Kette
auf $S = \{E, A, B, C\}$ (Zustände nach EABC-Klassifikation). Sei $\pi$ die
stationäre Verteilung ($\pi \mathcal{L} = \pi$).

\begin{definition}[Ausnahme-Trajektorie]
Eine Trajektorie $(T^k(n))_{k\geq 0}$ heißt \emph{ausnahme-typisch} für den
Zustand $s \in S$, falls
\[
  \limsup_{n\to\infty}\frac{1}{n}\sum_{k=0}^{n-1}\mathbf{1}[T^k(n) \equiv s]
  > \pi(s) + \varepsilon
\]
für ein $\varepsilon > 0$.
\end{definition}

\begin{tcolorbox}[kingman, title={Was Birkhoff/RPF über Ausnahmen sagt}]
Unter der stationären Verteilung $\pi$ gilt für $\pi$-fast alle Startzustände:
\[
  \frac{1}{n}\sum_{k=0}^{n-1}\phi(B_k(n)) \xrightarrow{n\to\infty} \Epi[\phi].
\]
Ein ``ausnahme-typisches'' $n$ müsste im Widerspruch zu Birkhoff stehen ---
\textbf{wenn $n$ ein typischer Startpunkt ist}. Das ist der Kern:
der Birkhoff-Ergodensatz gilt nur für $\mu_\pi$-fast alle Starts,
und $n \in \N$ könnte auf einer invarianten $\mu_\pi$-Nullmenge liegen.
\end{tcolorbox}

\subsection{Was Schritt A leistet und was nicht}

\begin{tcolorbox}[warnung, title={Grenzen von Schritt A}]
\begin{enumerate}[label=(\arabic*), nosep]
  \item Der Birkhoff-Ergodensatz sagt: Für $\mu_\pi$-fast alle $n$ ist
    eine dauerhaft ausnahme-typische Trajektorie \emph{unmöglich}.
  \item Aber: Gilt das auch für spezifische $n \in \N$? Nur wenn $\N$
    keine $\mu_\pi$-invariante Nullmenge bildet.
  \item Das ist eine zusätzliche Voraussetzung, die --- in allgemeiner
    Form --- zur Ergodizität von $T$ äquivalent ist.
\end{enumerate}
\textbf{Schritt A reduziert das Problem, löst es aber nicht.} Er zeigt,
dass ``zufällige'' Trajektories keine ewigen Ausnahmen sind. Für
konkrete Zahlen braucht man Schritt B oder C.
\end{tcolorbox}

% ====================================================
\section{Schritt B: mod-12-Primitivität und exponentielle Mischung}
\label{sec:schrittB}
% ====================================================

\subsection{Die Übergangsmatrix der mod-12-Markov-Kette}

Die mod-12-Markov-Kette auf $S = \{E, A, B, C\}$ hat Transfer-Operator
$\mathcal{L}$ mit Übergangsmatrix $M$. Numerisch gilt:

\begin{center}
\renewcommand{\arraystretch}{1.4}
\begin{tabular}{c|cccc}
\toprule
$M_{ij} = P(B_{n+1} = j \mid B_n = i)$ & E & A & B & C \\
\midrule
E & $3/8$ & $1/8$ & $3/8$ & $1/8$ \\
A & $3/8$ & $1/8$ & $3/8$ & $1/8$ \\
B & $3/8$ & $1/8$ & $3/8$ & $1/8$ \\
C & $3/8$ & $1/8$ & $3/8$ & $1/8$ \\
\bottomrule
\end{tabular}
\end{center}

\begin{bemerkung}
Die tatsächliche Übergangsmatrix hängt von den genauen Wahrscheinlichkeiten
ab, mit denen Blöcke auf Nachfolger des jeweiligen Typs zeigen. Die obige
Darstellung ist eine Vereinfachung; entscheidend ist, dass alle Einträge
\emph{positiv} sind (Primitivität), was nachfolgend angenommen wird.
\end{bemerkung}

\subsection{Irreduzibilität und Aperiodizität (Primitivität)}

\begin{tcolorbox}[rpf, title={Primitivität des Transfer-Operators (numerisch belegt)}]
Die mod-12-Markov-Kette auf $\{E, A, B, C\}$ ist:
\begin{itemize}[nosep]
  \item \textbf{Irreduzibel:} Jeder Zustand erreicht jeden anderen Zustand
    mit positiver Wahrscheinlichkeit (numerisch: alle $M^k$-Einträge
    für $k \leq 4$ positiv).
  \item \textbf{Aperiodisch:} Der $\gcd$ aller Rückkehrzeiten an jedem
    Zustand ist $1$ (Selbstschleifen durch E $\to$ E-Übergänge).
  \item \textbf{Primitiv} (= irreduzibel + aperiodisch): $M^k \to \mathbf{1}\pi^T$
    (Konvergenz zur stationären Verteilung) für $k \to \infty$.
\end{itemize}
\end{tcolorbox}

\subsection{Das exponentielle Mischungsargument}

\begin{theorem}[Mischungsrate des RPF-Operators]
\label{thm:mischung}
Sei $M$ primitiv mit stationärer Verteilung $\pi$ und zweitem Eigenwert
$|\lambda_2| < 1$. Dann gilt für alle Startzustände $s \in S$ und
alle $n \geq 0$:
\[
  \bigl\|M^n e_s - \pi\bigr\|_{\TV} \leq C\,|\lambda_2|^n,
\]
wobei $C > 0$ eine von $s$ abhängige Konstante ist.
\end{theorem}

\textbf{Numerisch:} $|\lambda_2| \approx 0{,}35$ (aus der Spektralanalyse
des mod-12-Operators).

\begin{tcolorbox}[ergebnis, title={Schlüsseleigenschaft: Mindestbesuchswahrscheinlichkeit}]
Sei $p := \min_{s,t \in S} M_{st} > 0$ die kleinste Übergangswahrscheinlichkeit.
Für eine Trajektorie der Länge $n$ gilt:
\[
  P\bigl(\text{Trajektorie besucht Zustand } E \text{ in den ersten }
  n \text{ Blöcken nie}\bigr) \leq (1-p)^n.
\]
Allgemeiner: Die Wahrscheinlichkeit, dass die Trajektorie $\lfloor n/12 \rfloor$
Schritte ohne ``gute'' Kontraktion verbringt, fällt exponentiell:
\[
  P\bigl(\text{keine gute Kontraktion in } \lfloor n/12 \rfloor \text{ Schritten}\bigr)
  \leq (1-p)^{\lfloor n/12 \rfloor}.
\]
\end{tcolorbox}

\subsection{Was Schritt B für das Uniformitätsproblem leistet}

\begin{tcolorbox}[kernidee, title={Schritt B: Stärken und Grenzen}]
\textbf{Stärken:}
\begin{enumerate}[label=(\arabic*), nosep]
  \item Die Mischungsrate ist \emph{explizit quantifizierbar} via $|\lambda_2|$.
  \item Für jeden konkreten Startzustand $s \in S$ (mod-12-Klasse) gibt
    der RPF-Satz eine \emph{gleichmäßige} Schranke für die Mischungszeit.
  \item Im Gegensatz zu Birkhoff braucht RPF keine Maßerhaltung ---
    nur Irreduzibilität und Aperiodizität der Kette.
\end{enumerate}
\textbf{Grenzen:}
\begin{enumerate}[label=(\arabic*), nosep]
  \item Der RPF-Satz für die mod-12-Kette garantiert Mischung
    auf Blockebene --- nicht für die ursprüngliche Abbildung $T$ auf $\N$.
  \item Ein konkretes $n \in \N$ könnte ``pathologische'' Blockfolgen
    erzeugen, die nicht der RPF-stationären Verteilung folgen, wenn
    die Blockabbildung $T^{(1)}$ nicht messbar-typisch wirkt.
  \item Die Aussage ``alle $n \in \N$ sind typisch'' bleibt eine
    zusätzliche Annahme, die formal nicht folgt.
\end{enumerate}
\end{tcolorbox}

\begin{bemerkung}[Wert für Taos Ergebnis]
Schritt B liefert: Die Dichte der Ausnahmen kann höchstens $(1-p)^{N/12}$
betragen, was für $N \to \infty$ gegen $0$ fällt. Das ist konsistent
mit $\mu_H(E) = 0$ (Tao) und quantifiziert die Mischungsgeschwindigkeit.
Aber $E = \emptyset$ folgt nicht.
\end{bemerkung}

% ====================================================
\section{Schritt C: Quaternionen-Diophantik (kritische Analyse)}
\label{sec:schrittC}
% ====================================================

\subsection{Das Argument in präziser Form}

Die Quaternionen-Norm-Identität lautet: Für die Hurwitz-Quaternionen
$\H$ und $q = a + bi + cj + dk$ mit $N(q) = a^2 + b^2 + c^2 + d^2$
gilt bei einer natürlichen Einbettung $n \mapsto q_n$:
\[
  N(3q + 1) = 9N(q) + 6a + 1.
\]
Diese Identität ist via \texttt{ring} in Lean beweisbar (für die
Einbettung $a = n$, $b = c = d = 0$):
\[
  N(3n+1) = (3n+1)^2 = 9n^2 + 6n + 1 = 9N(n) + 6n + 1. \checkmark
\]

Das vorgeschlagene Widerspruchsargument:
\begin{enumerate}[label=(\arabic*)]
  \item Für jede \emph{ganzzahlige} Trajektorie nimmt $N(T^k(q))$
    nur ganzzahlige Werte an.
  \item ``Fraktale'' Ausnahmepfade müssten irrational-algebraische
    Normwerte erzeugen.
  \item Widerspruch zur Ganzzahligkeit $\Rightarrow$ Ausnahmen können
    keine gewöhnlichen Ganzzahlen sein.
\end{enumerate}

\subsection{Kritische Analyse: Ist der Widerspruch wasserdicht?}

\begin{tcolorbox}[zirkel]
\textbf{Schritt C ist mathematisch \emph{nicht} wasserdicht.}
Es liegt eine versteckte Zirkularität vor.

\medskip
\textbf{Warum der Widerspruch nicht eintreten kann:}

Die Collatz-Abbildung $T: \N_{\mathrm{odd}} \to \N_{\mathrm{odd}}$,
$T(n) = (3n+1)/2^{\vii(3n+1)}$, ist definitionsgemäß eine Selbstabbildung
der ungeraden natürlichen Zahlen. Für jedes $n \in \N$ gilt also:
\[
  T^k(n) \in \N \quad\text{für alle } k \geq 0,
\]
\emph{unabhängig davon}, ob die Bahn konvergiert oder nicht. Ein
hypothetisches $n \in E$ hat eine Collatz-Bahn, die aus lauter
positiven ungeraden ganzen Zahlen besteht --- entweder divergiert sie
oder ist periodic. In beiden Fällen gilt:
\[
  N(T^k(n)) = (T^k(n))^2 \in \N \quad\text{für alle } k \geq 0.
\]

\medskip
\textbf{Die versteckte Zirkularität:}
Der behauptete Schritt (2) --- ``fraktale Ausnahmepfade erzeugen
irrational-algebraische Normwerte'' --- trifft schlicht nicht zu.
Ausnahmepfade sind Folgen ganzer Zahlen; ihre Quaternionen-Normen
sind ganzzahlig. Es entsteht kein Widerspruch zur Ganzzahligkeit.

\medskip
\textbf{Die zirkuläre Annahme:} Um zu begründen, warum Ausnahmepfade
``irrational-algebraische'' Strukturen erzeugen müssten, bräuchte man
bereits zu wissen, dass sie sich \emph{anders} als Ganzzahlfolgen
verhalten --- was genau das ist, was bewiesen werden soll.

\medskip
\textbf{Was die Identität tatsächlich leistet:} Die Norm-Identität
$N(3q+1) = 9N(q)+6a+1$ ist eine \emph{reine Rechenidentität} in $\H$.
Sie liefert keine neue Constraint auf die Collatz-Dynamik.
\end{tcolorbox}

\subsection{Kann Schritt C gerettet werden?}

Eine abgeschwächte Version von Schritt C könnte argumentieren:

\begin{tcolorbox}[konj]
\textbf{Mögliche Rettungsstrategie:} Falls man die Collatz-Dynamik
in einem Nicht-Standard-Sinne in $\H$ einbettet (z.\,B.\ via
2-adische Quaternionen $\H_2 = \H \otimes_\Q \Q_2$), dann könnten
Ausnahmepfade tatsächlich ``nicht-kompakt'' bezüglich einer
$p$-adischen Norm werden. Aber:
\begin{enumerate}[label=(\arabic*), nosep]
  \item Diese Einbettung müsste erst konstruiert werden.
  \item Der Widerspruch müsste zeigen, dass die 2-adische Quaternionen-Norm
    für divergente Bahnen unbeschränkt wird --- was wiederum auf die
    Divergenz selbst hinausläuft.
  \item Es ist nicht klar, ob diese Einbettung mehr liefert als
    das direkte Argument über $\Ztwo^*$.
\end{enumerate}
\textit{Diese Strategie ist unbewiesen und möglicherweise zirkulär.}
\end{tcolorbox}

\subsection{Ehrliche Einordnung von Schritt C}

\begin{center}
\renewcommand{\arraystretch}{1.4}
\begin{tabular}{>{\raggedright}p{4.0cm}
                >{\centering}p{2.8cm}
                >{\raggedright\arraybackslash}p{5.5cm}}
\toprule
\textbf{Aspekt} & \textbf{Bewertung} & \textbf{Anmerkung} \\
\midrule
Norm-Identität (Lean-beweisbar)
  & \textcolor{bewiesengrün}{\bfseries Klar wahr}
  & Ring-Identität in $\H$ \\[2pt]
Ganzzahligkeit der Normen entlang $\N$-Trajektorien
  & \textcolor{bewiesengrün}{\bfseries Trivial wahr}
  & Gilt auch für Ausnahme-Bahnen \\[2pt]
Widerspruch für Ausnahme-$n$
  & \textcolor{lückenrot}{\bfseries Nicht gegeben}
  & Zirkularität: Ausnahmen bleiben ganzzahlig \\[2pt]
``Fraktale irrational-algebraische Werte''
  & \textcolor{lückenrot}{\bfseries Nicht begründet}
  & Nicht aus der Identität ableitbar \\[2pt]
Gesamtstärke des Arguments
  & \textcolor{lückenrot}{\bfseries Schwach}
  & Kein echter Beitrag zum Uniformitätsproblem \\
\bottomrule
\end{tabular}
\end{center}

% ====================================================
\section{Das vollständige Uniformitätsargument}
\label{sec:uniform}
% ====================================================

\subsection{Was aus A, B, C zusammen folgt}

Kombiniert man Schritte A, B, C, ergibt sich folgendes Bild:

\begin{tcolorbox}[kernidee, title={Gesamtbild: Was gezeigt werden kann}]
\begin{enumerate}[label=(\arabic*)]
  \item \textbf{Tao + Schritt B (mod-12-RPF):} Der Ausnahmeset $E$
    hat Dichte~$0$ --- sowohl logarithmisch (Tao) als auch mit
    expliziter Mischungsrate $(1-p)^{n/12}$.
  \item \textbf{Schritt A + B:} Ein $n \in E$ müsste eine Trajektorie haben,
    die systematisch ``schlechte'' Blockklassen ansteuert. Die
    Mischungsrate des RPF-Operators macht solche dauerhaften Ausnahmen
    exponentiell unwahrscheinlich --- \emph{für generische Starts}.
  \item \textbf{Schritt C:} Kein Beitrag (Zirkularität).
\end{enumerate}

\medskip
\textbf{Gesamtresultat:} Die Kombination aus Tao und mod-12-RPF zeigt,
dass jedes $n \in E$ eine hochgradig unnatürliche arithmetische Struktur
benötigt --- eine Trajektorie, die die stationäre Blockverteilung $\pi$
systematisch meidet. Das ist heuristisch sehr unwahrscheinlich, aber
\emph{nicht formell ausgeschlossen}.
\end{tcolorbox}

\subsection{Warum das Uniformitätsproblem offen bleibt}

\begin{tcolorbox}[lücke, title={Die unüberbrückte Lücke}]
Das Uniformitätsproblem lässt sich wie folgt präzisieren:

\medskip
\textbf{Gegeben:} Der Transfer-Operator ist primitiv, $|\lambda_2| < 1$,
$\Lambda \approx -0{,}830 < 0$.

\medskip
\textbf{Gesucht:} Ein Argument, warum jede Collatz-Trajektorie
$n \in \N$ die RPF-stationäre Verteilung $\pi$ erreich --- nicht nur
$\mu$-fast alle.

\medskip
\textbf{Das Problem:} Der Satz von Birkhoff/RPF liefert Konvergenz
für ``$\mu$-fast alle'' Startpunkte. Für $n \in \N$ (abzählbare
Nullmenge) braucht man entweder:
\begin{enumerate}[label=(\alph*), nosep]
  \item \textbf{Quantitative Gleichmäßigkeit:} Eine Schranke
    $\|M^n e_{B(m)} - \pi\|_{\TV} \leq C\lambda_2^n$ für
    \emph{alle} $m \in \N$ (nicht nur fast alle).
  \item \textbf{Ausschluss spezieller Nullmengen:} Zeigen, dass kein
    $m \in \N$ auf einer $T$-invarianten $\mu$-Nullmenge liegt.
\end{enumerate}
Für (a) ist die RPF-Schranke auf dem mod-12-Niveau hinreichend, aber
man braucht noch, dass die Blockabbildung $T^{(1)}$ alle
$m \in \N$ gleichförmig auf die Markov-Kette projiziert.
Für (b) ist die Frage äquivalent zur Ergodizität --- und damit
zur Collatz-Vermutung selbst.
\end{tcolorbox}

\subsection{Eine bedingte Aussage}

\begin{proposition}[Bedingte Uniformität, formal]
\label{prop:bedingt}
Angenommen:
\begin{enumerate}[label=(V\arabic*)]
  \item\label{V1} Die mod-12-Markov-Kette ist irreduzibel und aperiodisch
    (mit Spektrallücke $1 - |\lambda_2| > 0$).
  \item\label{V2} Die Blockabbildung $T^{(1)}: \N_{\mathrm{odd}} \to \N_{\mathrm{odd}}$
    ist gleichmäßig verteilt über die mod-12-Klassen in folgendem Sinne:
    Für alle $n \in \N$ und alle $K \in \N$ gilt
    \[
      \frac{1}{K}\sum_{k=0}^{K-1}\mathbf{1}[T^{(k)}(n) \equiv s \pmod{12}]
      \xrightarrow{K\to\infty} \pi(s).
    \]
\end{enumerate}
Dann gilt für alle $n \in \N$:
\[
  \frac{1}{K}\sum_{k=0}^{K-1}\log_2 F(B_k(n)) \xrightarrow{K\to\infty} \Lambda < 0,
\]
was (für $\Lambda < 0$) die Collatz-Vermutung für $n$ impliziert.
\end{proposition}

\begin{beweis}
Voraussetzung~\ref{V2} ist genau die Birkhoff-Konvergenz für $n$.
Mit $\Lambda < 0$ folgt $F_k(n) \to 0$ geometrisch, also
$T^{(K)}(n) \to 0$, also letztlich Konvergenz gegen $1$.
\end{beweis}

\begin{bemerkung}
Voraussetzung~\ref{V2} ist stark --- sie ist für ein \emph{einzelnes} $n$
nicht aus dem RPF-Satz allein ableitbar. Sie ist die Umformulierung
der Ergodizitätsforderung für dieses spezifische $n$.
\end{bemerkung}

% ====================================================
\section{Was bleibt offen? Bewertungsmatrix}
\label{sec:offen}
% ====================================================

\subsection{Tabelle der Argumente}

\begin{tcolorbox}[bewmat]
\begin{center}
\renewcommand{\arraystretch}{1.5}
\begin{tabular}{>{\raggedright}p{4.2cm}
                >{\centering}p{1.6cm}
                >{\centering}p{1.6cm}
                >{\raggedright\arraybackslash}p{4.8cm}}
\toprule
\textbf{Argument} & \textbf{Stärke} & \textbf{Lücke?} & \textbf{Was fehlt} \\
\midrule
Tao: $\mu_H(E) = 0$
  & \textcolor{bewiesengrün}{\bfseries Stark}
  & Ja
  & $\N$ ist $\mu_H$-Nullmenge \\[2pt]
RPF mod-12: exponentielle Mischung
  & \textcolor{bewiesengrün}{\bfseries Mittel}
  & Ja
  & Formaler Spektrallückenbeweis offen \\[2pt]
LTE-Worst-Cases + strukturelle Kompensation
  & \textcolor{bewiesengrün}{\bfseries Mittel}
  & Ja
  & Uniforme Zweiblockkompensation \\[2pt]
Schritt A: Keine ewigen Worst-Cases
  & \textcolor{warnorange}{\bfseries Schwach}
  & Ja
  & Setzt Ergodizität voraus \\[2pt]
Schritt B: Mindestbesuchsrate $p > 0$
  & \textcolor{warnorange}{\bfseries Mittel}
  & Ja
  & Einzelnes $n$ vs.\ fast alle \\[2pt]
Schritt C: Quaternionen-Diophantik
  & \textcolor{lückenrot}{\bfseries Sehr schwach}
  & Ja
  & Zirkulär (§\,4) \\[2pt]
\midrule
\textbf{Gesamtschau}
  & \textcolor{warnorange}{\bfseries Offen}
  & \textbf{Ja}
  & Uniformitätsproblem nicht gelöst \\
\bottomrule
\end{tabular}
\end{center}
\end{tcolorbox}

\subsection{Nächste substanzielle Schritte}

\begin{enumerate}[label=(\Roman*)]
  \item \textbf{Spektrallücke formell beweisen:} Die Irreduzibilität
    und Aperiodizität der mod-12-Markov-Kette ist \emph{endliche Kombinatorik}.
    Der Beweis, dass alle Zustände in $\{E,A,B,C\}$ kommunizieren
    und dass die Periodizität $= 1$ ist, sollte formalisierbar sein
    (z.\,B.\ in Lean mit \texttt{Finset}-Methoden).

  \item \textbf{Zweiblockkompensation gleichmäßig:} Aus
    \textit{collatz\_lte\_dichte.tex}, Prop.\,3.1 folgt, dass nach dem
    LTE-extremalen Block eine Kontraktionsphase einsetzt. Kann man dies
    \emph{uniform} über alle $(k,r)$ zeigen? Das würde das
    Dichte-Argument auf Zwei-Block-Ebene retten.

  \item \textbf{Von mod-12 zu konkreten $n$:} Finden einer expliziten
    Klasse $\mathcal{C} \subseteq \N$ mit $\N \setminus \mathcal{C}$
    endlich, für die V2 (gleichmäßige Blockverteilung) beweisbar ist.

  \item \textbf{Schritt C ersetzen:} Das Quaternionen-Argument sollte
    durch einen vollständig neuen Ansatz ersetzt werden, der nicht
    auf der trivialen Ganzzahligkeit der Norm aufbaut.
\end{enumerate}

\subsection{Ehrliche Abschlussbewertung}

\begin{tcolorbox}[warnung, title={Ehrliche Gesamtbilanz}]
\begin{itemize}[nosep]
  \item Das dreistufige Argument (A, B, C) \textbf{schließt die Lücke
    zwischen Tao und der vollen Conjecture nicht}.
  \item Schritt B (mod-12-RPF) ist der mathematisch substanziellste
    Beitrag: Er quantifiziert, wie schnell Ausnahmen ``verdrängt''
    werden, und gibt $E$ eine exponentielle Mischungsschranke.
  \item Schritt C ist \textbf{mathematisch nicht valide} (Zirkularität)
    und sollte nicht als Bestandteil eines Beweisversuchs präsentiert werden.
  \item Schritt A ist konzeptionell korrekt, aber konditionell auf
    Voraussetzungen, die zur Collatz-Vermutung äquivalent sind.
  \item \textbf{Der nächste echte Schritt} ist der formale Beweis der
    Spektrallücke des mod-12-Transfer-Operators --- endliche Kombinatorik,
    Lean-formalisierbar.
\end{itemize}
\end{tcolorbox}

\subsection{Referenzen}

\begin{itemize}
  \item T.\ Tao: \textit{Almost all orbits of the Collatz map attain almost
    bounded values}. Forum Math.\ Pi \textbf{10} (2022), e12.
  \item J.\ F.\ C.\ Kingman: \textit{The ergodic theory of subadditive stochastic
    processes}. J.\ Roy.\ Stat.\ Soc.\ B \textbf{30} (1968), 499--510.
  \item D.\ Ruelle: \textit{Thermodynamic Formalism}. Addison-Wesley, 1978.
  \item Ergänzende Dokumente: \textit{collatz\_kingman.tex},
    \textit{collatz\_lte\_dichte.tex}, \textit{collatz\_schlussartikel.tex}.
\end{itemize}

% ====================================================
\section{Das Binärinformations-Argument: Analyse und Einordnung}
\label{sec:binaer}
% ====================================================

\subsection{Das intuitive Argument}
\label{sec:binaer:intuitiv}

\begin{tcolorbox}[colback=orange!9, colframe=warnorange, fonttitle=\bfseries,
  title={Argument D: Endliche Binärinformation (vielversprechend, aber unvollständig)}]
Jede natürliche Zahl $x \in \N$ besitzt eine \textbf{endliche Binärdarstellung}
mit $\lfloor\log_2 x\rfloor + 1$ Bits. Das Argument lautet:

\medskip
\textit{``Sobald die Trajektorie die ``Länge'' dieser Binärdarstellung abgelaufen hat,
bricht die lokale Speicherfähigkeit des Startwerts zusammen. Die exponentielle
Mischungsrate $(1-p)^{\lfloor n/12\rfloor}$ frisst die endliche Information
in endlicher Zeit auf.''}

\medskip
\textbf{Intuitive Stärke:} Der Startwert $x$ trägt nur $\lfloor\log_2 x\rfloor + 1$
Bits Information. Nach hinreichend vielen Collatz-Schritten --- mit exponentiell
mischender mod-12-Dynamik --- sollte diese endliche Information gegen die
stationäre Verteilung $\pi$ ``eingeebnet'' sein: kein Startwert $x$ kann
dauerhaft bestimmte mod-12-Klassen bevorzugen.

\medskip
\textbf{Status:} Das Argument ist intuitiv ansprechend und greift die richtige
Idee auf (endliche Information des Starts), aber enthält --- wie in \S\,\ref{sec:binaer:zirkel}
gezeigt --- eine subtile Zirkularität.
\end{tcolorbox}

\subsection{Die präzise Zirkularitätsprüfung}
\label{sec:binaer:zirkel}

\begin{tcolorbox}[lücke, title={Das zentrale Problem: Bitlänge wächst mit der Trajektorie}]
Der Collatz-Schritt $n \mapsto 3n+1$ kann die Bitlänge \emph{vergrößern}:
\[
  n \text{ hat } \lfloor\log_2 n\rfloor + 1 \text{ Bits},
  \qquad
  3n+1 \text{ hat ca.\ } \lfloor\log_2 n\rfloor + 2 \text{ Bits.}
\]
Präziser gilt:
\[
  \lfloor\log_2(3n+1)\rfloor \approx \lfloor\log_2 n\rfloor + \log_2 3
  \approx \lfloor\log_2 n\rfloor + 1{,}58.
\]
Die ``endliche Information'' des Startwerts $x$ wird durch die wachsenden
Folgeglieder \emph{erweitert}, nicht ``aufgefressen''. Nach $k$ ungeraden
Collatz-Schritten kann die Bitlänge bis auf
$\lfloor\log_2 x\rfloor + k \cdot \log_2 3 \approx \lfloor\log_2 x\rfloor + 1{,}58\,k$
angewachsen sein.
\end{tcolorbox}

\begin{proposition}[Äquivalenz mit Nicht-Divergenz]
\label{prop:bitlaenge}
Die Aussage ``Die Bitlänge der Trajektorie $(T^k(n))_{k\geq 0}$ bleibt nach
oben beschränkt'' ist \textbf{äquivalent zur Collatz-Vermutung für $n$} (im Sinne
der Nicht-Divergenz).
\end{proposition}

\begin{beweis}
$(\Rightarrow)$: Ist die Bitlänge beschränkt, so ist die Trajektorie beschränkt,
also periodisch. Da die einzige bekannte Periode die $1 \to 4 \to 2 \to 1$-Schleife
ist, folgt Konvergenz gegen $1$.

$(\Leftarrow)$: Konvergiert die Trajektorie gegen $1$, so ist sie letztlich beschränkt,
also auch die Bitlänge.
\end{beweis}

\begin{tcolorbox}[zirkel]
\textbf{Die Zirkularität von Argument D:} Man darf nicht annehmen, dass die
Trajektorie beschränkt bleibt --- und damit nicht, dass ihre Bitlänge durch
$\lfloor\log_2 x\rfloor + \mathrm{const}$ beschränkt ist. Genau diese
Annahme ist das, was bewiesen werden soll.

\medskip
\textbf{Vergleich mit Schritt C:} Argument D enthält dieselbe versteckte
Zirkularität wie Schritt C (\S\,\ref{sec:schrittC}): Beide setzen implizit
voraus, dass die Trajektorie sich ``kontrolliert'' verhält, um daraus die
Kontrolle abzuleiten.
\end{tcolorbox}

\subsection{Was korrekt bleibt: 2-adische Endlichkeit vs.\ Konvergenz}
\label{sec:binaer:adisch}

Trotz der Zirkularität enthält Argument D einen richtigen Kern, der präzisiert
werden kann:

\begin{tcolorbox}[adisch, title={Korrekte Aussagen zur 2-adischen Struktur}]
\begin{enumerate}[label=(\arabic*), nosep]
  \item \textbf{Endlichkeit in $\N$ vs.\ $\Ztwo$:} Für $n \in \N$ ist die
    Binärdarstellung endlich. Für allgemeines $x \in \Ztwo$ ist die
    2-adische Darstellung unendlich. Elemente $\N \subset \Ztwo^*$ sind
    \emph{isoliert} in der 2-adischen Topologie.

  \item \textbf{Fixpunkte auf $\Ztwo$:} Die Collatz-Abbildung auf $\Ztwo$
    besitzt Fixpunkte, z.\,B.\ $x = -1$ (da $3\cdot(-1)+1 = -2 \to -1$).
    Elemente von $\N$ liegen weit von diesen $p$-adischen Fixpunkten entfernt
    --- aber ob sie zu $1$ konvergieren, folgt daraus nicht.

  \item \textbf{Tao's 2-adische Methode:} Taos Beweis nutzt gerade die
    \emph{unendliche} 2-adische Entwicklung für die Fourier-Analyse auf $\Ztwo^*$.
    Die Endlichkeit der Darstellung von $n \in \N$ hilft dort nicht direkt.
\end{enumerate}
\end{tcolorbox}

\begin{bemerkung}[Warum Endlichkeit nicht genügt]
Sei $x \in \N$ mit Binärdarstellung der Länge $L = \lfloor\log_2 x\rfloor + 1$.
Die Tatsache $L < \infty$ unterscheidet $x$ topologisch von generischen
$\Ztwo$-Elementen, liefert aber keine arithmetische Schranke für die Dynamik.
Die Collatz-Abbildung ist nicht isometrisch bzgl.\ der 2-adischen Norm:
$|T(n)|_2$ kann kleiner oder größer als $|n|_2$ sein.
\end{bemerkung}

\subsection{Das stärkere, bedingte Argument}
\label{sec:binaer:bedingt}

Eine \emph{korrekte} und substanzielle Version des Informations-Arguments ist
die folgende bedingte Aussage, die RPF für die Konditionierung nutzt:

\begin{tcolorbox}[kernidee, title={Das bedingte Informations-Argument (korrekt)}]
Sei $n \in \N$ gegeben. Definiere $K(n)$ als den Index, nach dem die
mod-12-Klasse von $T^k(n)$ ``gut gemischt'' ist --- d.\,h.\ den kleinsten
$K$, so dass:
\[
  \bigl\|M^K e_{B(n)} - \pi\bigr\|_{\TV} \leq \varepsilon
\]
für ein vorgewähltes $\varepsilon > 0$. Dann gilt:
\[
  K(n) \leq \frac{\log(C \cdot \varepsilon^{-1})}{\log(|\lambda_2|^{-1})},
\]
wobei $C > 0$ von der Startklasse $B(n) \in \{E,A,B,C\}$ abhängt (nur 4 Fälle).
Dieser Wert $K(n)$ ist \textbf{endlich und explizit berechenbar} --- unabhängig
von der Größe von $n$.

\medskip
\textbf{Was das bedeutet:} Nach $K(n)$ Blockschritten ist die mod-12-Klasse
gut gemischt, \emph{falls die Trajektorie bis dahin nicht divergiert}. Das ist:
\begin{enumerate}[label=(\alph*), nosep]
  \item Kein Beweis der Nicht-Divergenz (conditio sine qua non).
  \item Aber: eine \emph{quantitative} Verstärkung gegenüber Schritt B.
    Die Mischungszeit hängt \emph{nicht} von $n$ ab, sondern nur von der
    mod-12-Klasse des Startwerts.
  \item Damit: Die ``endliche Binärinformation'' von $n$ verliert ihren
    Einfluss auf die Klassen-Dynamik nach $K(n)$ Schritten --- \emph{bedingt}
    auf Nicht-Divergenz.
\end{enumerate}
\end{tcolorbox}

\begin{proposition}[Klassen-Vergessen nach $K(n)$ Schritten]
\label{prop:vergessen}
Sei $p := \min_{s,t \in S} M_{st} > 0$ die kleinste Übergangswahrscheinlichkeit
der mod-12-Kette, und sei $|\lambda_2| < 1$ die Spektrallücke. Dann gibt es
$C > 0$ (unabhängig von $n$), so dass für alle $n \in \N$ und alle $k \geq 0$:
\[
  \bigl\|M^k e_{B(n)} - \pi\bigr\|_{\TV} \leq C \cdot |\lambda_2|^k.
\]
Insbesondere: Nach
\[
  K_\varepsilon := \left\lceil\frac{\log(C/\varepsilon)}{\log(1/|\lambda_2|)}\right\rceil
  \quad \text{Schritten}
\]
ist $\|M^{K_\varepsilon} e_{B(n)} - \pi\|_{\TV} \leq \varepsilon$
für \textbf{jeden} Startpunkt $n \in \N$ (da $K_\varepsilon$ nicht von $n$ abhängt).
\end{proposition}

\begin{beweis}
Direkt aus Theorem~\ref{thm:mischung} mit $C = \max_{s \in S} C_s$ über
die endlich vielen Startklassen $s \in \{E, A, B, C\}$.
\end{beweis}

\begin{tcolorbox}[bewmat, title={Vergleichende Einordnung der vier Argumente}]
\begin{center}
\renewcommand{\arraystretch}{1.5}
\begin{tabular}{>{\raggedright}p{2.0cm}
                >{\centering}p{2.2cm}
                >{\centering}p{1.8cm}
                >{\raggedright\arraybackslash}p{5.2cm}}
\toprule
\textbf{Argument} & \textbf{Stärke} & \textbf{Zirkulär?} & \textbf{Kernaussage} \\
\midrule
Schritt A (Ergodentheorie)
  & \textcolor{warnorange}{\bfseries Bedingt}
  & Ja (implizit)
  & Keine ewigen Ausnahmen --- setzt Ergodizität voraus \\[2pt]
Schritt B (RPF mod-12)
  & \textcolor{bewiesengrün}{\bfseries Quantitativ}
  & Nein
  & Mischung in $K_\varepsilon$ Schritten, unabhängig von $n$ \\[2pt]
Schritt C (Quaternionen)
  & \textcolor{lückenrot}{\bfseries Schwach}
  & \textbf{Ja}
  & Kein echter Beitrag (§\,\ref{sec:schrittC}) \\[2pt]
Arg.\ D (Binärinfo)
  & \textcolor{warnorange}{\bfseries Bedingt}
  & \textbf{Ja}
  & Korrekt nur bedingt auf Nicht-Divergenz \\
\bottomrule
\end{tabular}
\end{center}
\end{tcolorbox}

\begin{tcolorbox}[warnung, title={Fazit zu Argument D}]
Das Binärinformations-Argument ist \textbf{intuitiv wertvoll}, weil es auf
den richtigen Mechanismus zeigt: Die endliche ``Gedächtnistiefe'' eines
Startwerts $n \in \N$ sollte nach hinreichend vielen Schritten gegen die
stationäre Verteilung eingeebnet werden. Der präzise Inhalt dieser Idee
ist Proposition~\ref{prop:vergessen} (keine Abhängigkeit von $n$ in $K_\varepsilon$).

Zur Schließung der Uniformitätslücke fehlt jedoch der entscheidende Schritt:
der Nachweis, dass die Trajektorie \emph{überhaupt} $K_\varepsilon$ Schritte
macht, ohne zu divergieren. Das ist zirkulär.

\textbf{Nächster echter Schritt:} Kombiniere Prop.~\ref{prop:vergessen}
mit einer \emph{Divergenzausschluss}-Schranke --- z.\,B.\ über 2-adische
Kontraktionsargumente oder eine positive untere Schranke für die
Konvergenzwahrscheinlichkeit in den ersten $K_\varepsilon$ Schritten.
\end{tcolorbox}

\end{document}
